\subsection{Overview}

Agricultural subsurface (tile) drainage is one of the most important
anthropogenic modifiers of the field-to-stream nitrogen and phosphorus
pathway, and a large, geographically diverse body of field, watershed,
and modelling evidence now supports a fairly consolidated picture of how
it controls nutrient export. The central, repeatedly confirmed finding
is that subsurface drainage is the dominant hydrologic and nitrogen loss
pathway in artificially drained cropland: in intensively monitored
Norwegian catchments, subsurface drainage carried $76$--$91\%$ of runoff
and $87$--$95\%$ of total-nitrogen loss, with drainage total-N
concentrations $2$--$4$ times higher than in surface runoff and over
$90\%$ of TN losses occurring through subsurface drainage
\parencite{Bechmann2014}. Nitrate-N is the predominant exported species,
and its transport is tightly coupled to discharge, so that hydrology and
fertilizer intensity jointly govern loads \parencite{Raymond2012}. Across
the literature the headline quantities recur: field- and watershed-scale
nitrate-N yields commonly reach or exceed $20$--$60$~kg~ha$^{-1}$~yr$^{-1}$
in maize--soybean systems \parencite{Gentry2009,Gentry2014,Dattamudi2020,Tomer2008},
concentrations frequently approach or exceed the $10$~mg~L$^{-1}$
drinking-water limit \parencite{Dattamudi2020,Gordon2021}, and phosphorus,
though smaller in mass, moves through both tile and surface pathways at
concentrations well above eutrophication thresholds \parencite{Smith2008,Tomer2008}.

The evidence also converges on a suite of mitigation strategies with
quantified, mechanistically coherent effects. Drainage water management
(controlled drainage, subirrigation, and drainage water recycling)
reliably reduces drainage volume and nitrate load, typically in
proportion to the flow reduction, while often maintaining or increasing
yield \parencite{Skaggs2010,Youssef2018,Wu2025,Madramootoo2024}.
Edge-of-field and beyond-the-field practices, constructed wetlands,
denitrification bioreactors and walls, saturated buffers, two-stage
ditches, ditch weirs, riparian zones, and filter strips, remove nitrate
mainly by denitrification and reduce phosphorus and sediment with more
variable efficiency \parencite{Tanner2005,Jaynes2008,Hodaj2017,Strock2025b,JACOBS1985}.
In-field practices, principally cover crops and nutrient management, cut
tile nitrate loss substantially at the tile scale but show attenuated,
scale-dependent signals at the watershed outlet
\parencite{Hanrahan2018,Speir2022,Gentry2014}. This section synthesizes
these consolidated findings, organized around the magnitude and pathways
of loss, the hydrologic and land-management controls that drive them, and
the drainage-management, edge-of-field, and watershed-scale strategies
shown to abate them, together with the salinity, trace-element, and
pesticide dimensions of drainage water quality.

\subsection{Magnitude and Pathways of Nitrogen Loss}

\subsubsection{Subsurface drainage as the dominant nitrogen pathway}

Multiple independent studies establish that in tile-drained landscapes
the subsurface pathway carries the overwhelming majority of nitrogen.
In Norway, over $90\%$ of TN losses occurred through subsurface drainage,
concentrations in drainage were $2$--$4$ times those in surface runoff,
average agricultural-area TN concentrations ranged $3$--$19$~mg~L$^{-1}$,
and nitrate-N constituted $71$--$85\%$ of TN except at one site
\parencite{Bechmann2014}. New Zealand mole-pipe systems captured about
$62\%$ of surplus rainfall as drainage and routed roughly $80\%$ of
soluble N discharged in combined flows, while surface runoff dominated
P, sediment, and \emph{E. coli} transfer \parencite{Monaghan2016}. In
east-central Illinois, tile drainage losses were a major N source to
surface waters, with tile N loads of $22.7$ and $59.9$~kg~ha$^{-1}$ in
successive years and flow-weighted mean NO$_3$-N rising to $19.2$~mg~L$^{-1}$
in the wetter year \parencite{Gentry2009}. Watershed nitrate export in
Iowa exceeded $20$~kg~ha$^{-1}$~yr$^{-1}$ and tracked baseflow used as a
surrogate for tile flow \parencite{Tomer2008}, and long-term Illinois
budgets showed subsurface stations losing $13$--$30$~kg~ha$^{-1}$~yr$^{-1}$
of nitrate-N (equal to $16$--$26\%$ of applied N) with more than $70\%$ of
subsurface samples exceeding the $10$~mg~L$^{-1}$ MCL, while surface runoff
was negligible in the flat terrain \parencite{Dattamudi2020}.

\subsubsection{Loads, concentrations, and their variability}

Reported magnitudes span a wide but coherent range. New Zealand tile
drains exported $14$--$109$~kg~ha$^{-1}$~yr$^{-1}$ of TN, of which
$58$--$90\%$ was nitrate-N, with annual drainage yields of $193$--$853$~mm
\parencite{Tanner2011}. Riverine and tile nitrate yields in Illinois
averaged around $30$--$32$~kg~N~ha$^{-1}$~yr$^{-1}$ and ranged up to
$57$--$59$~kg~ha$^{-1}$~yr$^{-1}$ \parencite{Gentry2014}, and Indiana ditch
networks showed NO$_3$-N loads of $0.05$--$19$~kg~ha$^{-1}$ with
concentrations peaking above $16$--$20$~mg~L$^{-1}$ in wet springs
\parencite{Smith2008}. Across 29 Iowa watersheds, five-year moving-average
flow-weighted nitrate concentrations ranged $0.6$--$15.5$~mg~L$^{-1}$,
loads spanned $0.07$--$58\times10^{6}$~kg~yr$^{-1}$, and yields ranged
$3.6$--$78.7$~kg~ha$^{-1}$; crucially, interannual variability in load
greatly exceeded that of concentration, so flow-weighted concentration
served as a more robust early indicator of management change than load or
yield \parencite{Danalatos2022}. In a Czech watershed, flow-weighted
drainage nitrate rose with risk class from $9.6$ to $17.8$~mg~L$^{-1}$,
corresponding to drainage nitrate contributions of $796$--$2290$~kg~N~km$^{-2}$~yr$^{-1}$
\parencite{Kvítek2023}. German catchment analysis documented a decline in
drainage N discharge from $33$ to $10$~kg~N~ha$^{-1}$~yr$^{-1}$ (1480 to
454~t~yr$^{-1}$) between the 1980s and 1990s, driven mainly by reduced N
surplus, with the largest area-specific losses on sandy and loess
substrates \parencite{Hirt2005b}.

\subsubsection{Fertilizer input and hydrology as joint drivers}

At regional scale, fertilizer input and watershed hydrology together
regulate export: about $34\%$ of applied fertilizer N was exported to
streams in the Mississippi basin, wetter watersheds exported a larger
fraction, and tile drainage was identified as a hydrologic modifier that
increases throughput and reduces residence time, thereby raising nitrate
loss to the Gulf of Mexico \parencite{Raymond2012}. Direct comparison of
fertilized and unfertilized Minnesota fields confirmed the causal link:
nitrate load from the unfertilized field was $58\%$ of that from the
fertilized field \parencite{Ghane2016b}. Watershed N budgets showed
fertilizer as the largest input and grain as the largest output, with
mass balances swinging from positive to strongly negative in wet years
that depleted soil N \parencite{Gentry2009}, and plant uptake dominating
outputs ($89$--$96\%$) while roughly $80\%$ of residual soil-profile N was
lost to runoff and leaching \parencite{Dattamudi2020}. Multivariate models
that combined current and prior-year water yield, prior corn yield, and
fall fertilizer sales explained $96\%$ of annual nitrate-yield variation,
and shifting fertilizer application to spring was estimated to reduce
nitrate yield by $17$--$20\%$, detectable in watersheds as large as
$481$~km$^2$ \parencite{Gentry2014}. Principal-component and regression
analyses similarly identified the hydrologic component, together with
flow, N application, and temperature, as the dominant controls on nitrate
load \parencite{Li2015}.

\subsection{Phosphorus and Sediment Loss}

While nitrogen dominates the mass balance, phosphorus and sediment loss
through drainage systems is well documented and follows distinct
controls. Total P loads are typically small ($0.4$--$0.7$~kg~ha$^{-1}$~yr$^{-1}$)
yet concentrations commonly exceed $0.1$~mg~L$^{-1}$ and respond to storm
runoff more than nitrate \parencite{Tomer2008}; in Ohio tiles, mean annual
DRP and TP loads were $0.46$ and $0.57$~kg~ha$^{-1}$, with DRP and TP
concentrations higher in the growing season \parencite{King2016}. Surface
inlets exert a distinct effect on drainage nutrient transport: site-years
with surface inlets had greater median TP load ($0.40$ vs.
$0.21$~kg~ha$^{-1}$) but smaller dissolved-N load ($3.3$ vs.
$23.0$~kg~ha$^{-1}$), and a logistic model using TP and DN loads correctly
classified $89\%$ of site-years, offering a tool to infer upstream
practices \parencite{Flores2021}. Phosphorus is delivered through both
surface and subsurface routes: in flat New Zealand pastures surface runoff
carried $54\%$ of both filterable-reactive and total P
\parencite{Monaghan2016}, whereas Swedish diagnosis identified subsurface
drainage itself as the single largest probable cause of P losses, with
stream proximity second \parencite{Djodjic2002}. Legacy and macropore
processes matter: winter cover-crop tiles showed unexpectedly high TP
during rain and rain-on-snow events, with concurrent TP/TOC increases and
nitrate declines indicating bypass flow through macropores
\parencite{McNeill2026}, and Irish ditch sediments accumulated legacy P
that made farmyard-connection and outlet ditches the highest-risk P
pathways \parencite{Moloney2020,Opoku2024}. Watershed export
also concentrates in spring snowmelt "hot moments," with identifiable
spatial hot spots for DRP and TP distinct from nitrate hot spots
\parencite{Irvine2019}.

\subsection{Hydrologic and Climatic Controls}

\subsubsection{Discharge, high flows, and seasonality}

Discharge is the master variable for both N and P export. Tile discharge
was the main driver of soluble reactive phosphorus load in Indiana
watersheds, and subsurface flow--nitrate load relationships were strongly
positive ($R^2 = 0.84$--$0.86$) \parencite{Trentman2020,Dattamudi2020}.
More than $70\%$ of watershed nitrate export occurred during high-flow
events and moist conditions \parencite{Hanrahan2018}, and high flows were
identified as the major control on contaminant transport, implying that
targeting high flows could reduce loads \parencite{Ghane2016b}. Discharge
is strongly seasonal: in Ohio, $81\%$ of tile discharge occurred in the
nongrowing season, so nitrate and TN loads were greater then despite
higher growing-season concentrations under corn \parencite{King2016}.
Controlled-drainage studies similarly show winter and spring dominance of
loss, with most nitrogen loss between April and May
\parencite{Strock2025b,Grenon2023}. Nutrient concentrations frequently
peak after soil thaw and in spring following planting, especially in
wetter years \parencite{Magner2004,Smith2008,Bechmann2014}. Antecedent
moisture modulates these dynamics: wetter antecedent conditions increased
tile discharge and SRP load, while discharge declined as days since the
last $10$ or $20$~mm event increased \parencite{Trentman2020}.

\subsubsection{Climate change projections}

Modelling indicates that intensifying hydrology will increase drainage
throughput. In an Indiana watershed, annual subsurface drain flow was
projected to increase by up to $70\%$ by century's end (up to $67\%$ under
RCP~8.5), with surface runoff rising up to $131\%$; nitrate loads in tile
drains were projected to fall $25$--$75\%$ by mid-century while surface-runoff
nitrate rose $50$--$100\%$, and soluble-P yield via drains declined
$30$--$60\%$ \parencite{Mehan2019}. Ensemble analysis for an Ohio watershed
found no consensus on the direction of future nutrient change, with
climate-model choice dominating uncertainty for discharge, tile discharge,
ET, and TN, and hydrologic-model choice dominating for runoff and
phosphorus; on average tile discharge rose $18\%$ while TP fell $7\%$ and
TN changed little \parencite{Kujawa2020}. Long-term Latvian and Lithuanian
records already show rising precipitation and temperature and shifting
seasonal outflow, with significant winter increases and spring decreases
in drainage \parencite{Lagzdins2015,Povilaitis2018}. Increasing
precipitation and discharge are expected to shorten N residence time and
exacerbate export \parencite{Raymond2012}.

\subsection{In-Field Management: Cover Crops, Rotation, and Nutrient Management}

\subsubsection{Cover crops at the tile scale}

Cover crops consistently reduce nutrient loss where measured directly at
the tile. In Indiana, tiles draining cover-cropped fields had
$69$--$90\%$ lower seasonal nitrate loads in winter and spring
\parencite{Hanrahan2018}, and independent monitoring found cover crops cut
tile discharge by $28$--$65\%$ and nitrate-N loss by $27$--$72\%$ annually,
with consistent winter--spring reductions and more variable SRP reductions
of $7$--$59\%$ \parencite{Speir2022}. Spring cover-crop biomass reduced
median tile discharge by $50$--$80\%$ and SRP load by up to $54\%$, with
the effect estimated at a $15\%$ reduction in SRP over the fallow period
\parencite{Trentman2020}. In Missouri terraced fields, a paired-watershed
design attributed a $30\%$ reduction in subsurface flow and $75\%$
reduction in TP loss to cover crops \parencite{SalcedaGonzalez2025}. Winter
cover choice matters: winter-wheat fields had the highest and mixed
cover-crop fields the lowest nongrowing-season N concentrations
\parencite{McNeill2026}. Economic analysis showed environmental
co-benefits recover a large share of costs, averaging $61\%$, with the
largest contribution from cover-crop biomass N mineralization
\parencite{Roth2018}.

\subsubsection{Scale attenuation and offsetting responses}

At the watershed outlet the signal is weaker and scale-dependent. Despite
strong tile-scale effects, annual, seasonal, and monthly watershed nitrate
export did not differ significantly between pre- and post-cover-crop years,
although average annual nitrate yield was $13\%$ lower, and monitored tiles
contributed only $\sim7\%$ of watershed export \parencite{Hanrahan2018}.
Watershed-scale reductions were most consistent in spring, and consistent
annual reductions likely require broader regional or targeted
implementation \parencite{Speir2022}. Some cover-crop and relay systems
showed no drainage or nitrate benefit and higher N surplus or N$_2$O
emissions: a winter camelina relay produced no change in cumulative
nitrate load but increased nitrous oxide emissions and N surplus
\parencite{Emmett2022}, and Ohio BACI trials found cover crops reduced
nitrate flow-weighted mean concentration but not nitrate load
\parencite{Askar2023}. Grassing of drainage recharge zones reduced median
nitrate by $26$--$32\%$ locally and monthly N load by up to $47\%$
\parencite{Fucík2015,Kvítek2023}.

\subsection{Drainage Water Management: Controlled Drainage and Subirrigation}

\subsubsection{Consolidated hydrologic and nitrogen effects}

Controlled drainage (CD) is among the best-established practices. A
synthesis of prior studies found CD reduced drainage outflow by $17\%$ to
over $80\%$ depending on soil, drainage intensity, control strategy, and
location, while raising water tables and increasing ET, seepage, and
surface runoff; for many studies nitrogen-loss reduction was about the
same percentage as the drainage-volume reduction, though the N effect
depended on denitrification and seepage pathways \parencite{Skaggs2010}. A
global meta-analysis of 84 publications, 52 sites, and 1788 paired
site-years confirmed that both CD and controlled drainage with
subirrigation (CDSI) reduce subsurface discharge and subsurface nutrient
losses while modestly increasing yield; CDSI reduced subsurface nitrate
and TN more than CD, whereas CD reduced subsurface TP more, and CDSI (but
not CD) significantly increased surface runoff and surface nitrate loss
\parencite{Wu2025}. Simulations across 48 U.S. Midwest locations found CD
reduced annual drainage and N losses everywhere while increasing runoff,
ET, and seepage, with efficacy strongly climate-dependent (largest
reductions at Memphis, smallest at Fargo) \parencite{Youssef2018}.

\subsubsection{Field evidence of nitrate reductions and yield effects}

Field studies quantify these effects concretely. In Quebec, controlled
drainage with subirrigation cut tile drainage volume by $24\%$, nitrate
concentration by $25\%$ ($10.6$ to $7.9$~mg~L$^{-1}$), and nitrate loss by
$43\%$ ($25.8$ to $14.6$~kg~ha$^{-1}$~yr$^{-1}$) \parencite{Ng2000}, and a
long-term Canadian study found CD reduced total nitrate loss by $\sim58\%$,
increased yield $3\%$, and cut drainage volume $17\%$ without elevated
greenhouse-gas emissions \parencite{Madramootoo2024}. In North Carolina,
nine years of CD reduced drainage outflow $33\%$ and nitrate-N export
$30\%$ (average $6.3$~kg~ha$^{-1}$~yr$^{-1}$), a reduction equal to the
extra N removed in grain as CD raised corn and soybean yields $10$--$11\%$
\parencite{Poole2018}. Controlled tile drainage during the growing season
reduced daily water flux $\sim60\%$, nitrate $51\%$, ammonium $58\%$, DRP
and TP $66\%$, and \emph{E. coli} $76\%$, while boosting corn and soybean
yields \parencite{Sunohara2016}, and the same practice imposed en masse
across a watershed significantly reduced growing-season stream fluxes of
water, N, and P, with physical abatement of tile flow as the overriding
mechanism \parencite{Sunohara2015}. Iowa DWM cut tile flow $21\%$ and
nitrate mass loss $29\%$ with an $8\%$ soybean-yield increase, though
flow-weighted concentration was not significantly changed
\parencite{Jaynes2012}. Missouri controlled drainage reduced NO$_3$-N
concentrations by $6.0$--$9.6$~mg~L$^{-1}$ and cumulative nitrate loads by
$69$--$99\%$ during growing seasons, but raised dissolved reactive P
concentrations \parencite{Nash2020}.

\subsubsection{Timing, trade-offs, and phosphorus concerns}

The timing of control governs efficacy: DRAINMOD analysis identified
$1$--$15$~March as the most effective window, when CD could reduce
subsurface outflow and nitrate-N losses by up to $100\%$ and retain up to
$22$~kg~NO$_3$-N~ha$^{-1}$, whereas April starts resembled free drainage
\parencite{Kesicka2023c}. Lithuanian CD reduced monthly outflow
$55$--$80\%$, inorganic-N load $42$--$77\%$, and TP load $34$--$72\%$
\parencite{Povilaitis2018}. Modelled shallow drainage outperformed CD for
nitrogen (reducing N loss $40\%$ vs. $26\%$ long-term) with minimal yield
penalty \parencite{Craft2018}. Trade-offs are recognized: CD can increase
surface DRP and TP losses that partly offset subsurface benefits
\parencite{King2022,Wu2025}, and controlled drainage--subirrigation was
predicted to raise runoff N ($+340\%$) and sediment N ($+118\%$) even while
cutting leaching $35\%$, for a net $17.8\%$ reduction in combined
non-point N \parencite{SHIRMOHAMMADI1995}. Greenhouse-gas responses are
mixed but generally not adverse: N$_2$O fluxes under CD were $49\%$ higher
than free drainage in three years and up to $43\%$ lower in three others
\parencite{Madramootoo2024}. Drainage-management alone often fails to lower
concentrations even when it lowers loads via volume reduction
\parencite{Grenon2023,Wesström2014}. Automatic CD/subirrigation
outperformed traditional CD, reducing N losses up to $27\%$, increasing
plant uptake, and saving $28$--$77\%$ of subirrigation water
\parencite{Abdelbaki2025}. Subirrigation lysimeters in Iran cut drainage
nitrate concentrations by $74$--$84\%$ and improved yield and water-use
efficiency \parencite{Noory2009}, and Swedish controlled drainage/subirrigation
reduced NO$_3$-N, TN, PO$_4$-P, and TP losses to about $40\%$ of
conventional while raising yield and N uptake \parencite{Wesström2014}.

\subsection{Drainage Water Recycling and Reuse}

Drainage water recycling (DWR), capturing and storing drainage for
supplemental irrigation, is documented to benefit both crop production and
downstream water quality. Reported yield gains reach $61\%$ for corn and
$31\%$ for soybean relative to free drainage, with modelled Ohio benefits
rising under future climate; Canadian studies documented N-loss reductions
of $40$--$70\%$ and P-loss reductions of $12$--$36\%$, and modelling
predicted downstream N-delivery reductions of $24$--$37\%$ and P reductions
of $21$--$39\%$ in Iowa and Indiana \parencite{Hay2021}. A North Carolina
reservoir reduced TN load $47\%$, TP load $30\%$, and sediment $87\%$,
driven by nitrate ($-53\%$) and orthophosphate ($-46\%$) concentration
reductions, with hydraulic retention time the key control, though it was a
net source of particulate P \parencite{Moursi2023}. An expanded DRAINMOD
DWR model showed reservoir inflow dominated by subsurface drainage
($67\%$), and a $4\%$-area reservoir optimized corn yield while reducing
nitrate loss $\sim20\%$ ($7.7$~kg~ha$^{-1}$~yr$^{-1}$)
\parencite{Moursi2022}. Multi-location synthesis of 53 site-years showed
DWR is more resilient than free drainage, protecting yield during drought
and excess precipitation, especially in the V9--R2 window, producing
$2.22$~Mg~ha$^{-1}$ higher yield overall and reducing yield variability
$22\%$ \parencite{Willison2021}. Adoption economics indicate crop insurance
is not a major barrier: real-options adoption thresholds shifted less than
$\sim20\%$ with insurance, and DWR and insurance manage different risks
\parencite{Sellars2022}. Regional DWR in Iran saved $67$~Mm$^3$ of water,
recycled $32$~Mm$^3$, and raised river flow while cutting drainage to
evaporation ponds $38\%$; optimal blending was $13$--$23\%$ recycled water,
and a $20\%$ scenario limited sugarcane yield loss to acceptable levels
\parencite{Kosari2024}. In California's San Joaquin Valley, integrated
on-farm drainage management reused all but $0.7\%$ of applied water, using
only $\sim6\%$ of project area for evaporation and reuse, though drainage
EC did not represent crop-water concentration effects
\parencite{Ayars2014}.

\subsection{Edge-of-Field and Beyond-the-Field Practices}

\subsubsection{Constructed wetlands}

Constructed wetlands treating drainage are extensively documented as
effective nitrate sinks, with denitrification the dominant mechanism.
Across drained French watersheds, nitrate removal averaged $\sim42\%$ (up
to $40$--$90\%$), dominated by denitrification ($93\%$) over plant uptake
($7\%$), and devoting $\sim1\%$ of contributing area to a wetland can
realistically achieve $50\%$ removal, with hydraulic residence time the
priority control \parencite{Tournebize2017}. Iowa wetlands reduced incoming
nitrate concentrations $49\%$ and loads $\sim61$~kg~day$^{-1}$, with
performance varying by hydrology ($23$--$65\%$) \parencite{Drake2018}; a
Minnesota wetland removed $67\%$ of tile nitrate over four years, mostly
via infiltration to shallow groundwater \parencite{Gordon2021}; and another
removed $68\%$ of nitrate and $\sim76\%$ of phosphorus, with plant harvest
the primary P-removal route \parencite{Lenhart2016}. Illinois wetlands
retained $36$--$37\%$ of nitrate and total N and over half of P fractions,
and $\sim450$~ha of wetland were estimated necessary to cut watershed N
loading $46\%$ \parencite{Kovacic2006}. New Zealand wetlands occupying
$1$--$2\%$ of catchment area removed $16$--$63\%$ of nitrate-N loads but
performed poorly for P, sometimes exporting more TP than received
\parencite{Tanner2005,Tanner2005b,Tanner2011,Tanner2003}. Wetland
performance is strongly a function of hydraulic loading rate, residence
time, temperature, and inlet concentration: a first-order areal model with
median $k_{20}=25$~m~yr$^{-1}$ and temperature coefficient $1.106$ best
represents whole-system behaviour, and no single-number design is adequate
\parencite{Kadlec2012,Kalcic2018}. A German wetland retained only $2.9\%$
of TN over four years, but retention reached $\sim80\%$ when loads were low
and residence time exceeded $20$~days at temperatures above $8\,^{\circ}$C
\parencite{Steidl2019}. Meta-analysis reported mean cropland-drainage
removals of $23.7\%$ TN and $30.3\%$ nitrate-N but negative mean removals
for ammonium and TP, underscoring variable P and ammonium behaviour
\parencite{Messer2021}.

\subsubsection{Bioreactors, denitrification walls, and stacked practices}

Woodchip-based denitrification is well supported. Denitrification walls
reduced flow-weighted nitrate concentrations by at least $55\%$ and cut
average annual nitrate mass loss $55\%$ over five years with no yield
penalty and no land taken out of production \parencite{Jaynes2008}, and a
19-year comparison found a rye cover crop and a woodchip denitrification
wall reduced N leaching $58$--$59\%$, with rye more effective in dry years
and the wall more effective in wet years \parencite{Rogovska2023}. Trench
backfill additives reduced tile nitrate by $52$--$78\%$ (woodchips best),
and among in-ditch filter media Polonite ranked highest for P removal
\parencite{Povilaitis2018}. A stacked DWM--woodchip bioreactor--saturated
riparian buffer system reduced nitrate-N export $78$--$88\%$ over five
seasons, sustaining performance even in the final year, though the system
could be a net source of suspended solids \parencite{Singh2024b}. Saturated
buffers were highlighted as offering the largest climate-mitigation
potential per unit area among conservation drainage practices, decreasing
nitrate leaching with little increase in N$_2$O and accruing SOC
\parencite{Hall2025}.

\subsubsection{Ditches, weirs, two-stage ditches, and riparian buffers}

In-channel and edge practices provide additional nitrate and P
interception. A ditch with a low-grade weir reduced discharge $51\%$,
nitrate concentration $22\%$, and nitrate load $58\%$ over seven years in a
cold climate without removing land from production \parencite{Strock2025b}.
A two-stage ditch significantly reduced TP, SRP, and TSS concentrations and
loads (TP load $-40\%$), and lowered nitrate concentrations, though not
nitrate loads \parencite{Hodaj2017}; two-stage ditches and wetlands are
regarded as promising beyond-the-field practices capturing both N and P
\parencite{Kalcic2018}. Riparian buffers are highly effective for nitrate:
in North Carolina, groundwater nitrate dropped from $7$--$8$~mg~L$^{-1}$ to
$<0.1$~mg~L$^{-1}$ within $16$--$47$~m of the drainageway, attributed to
denitrification rather than plant assimilation, and buffers $<16$~m
induced significant nitrate loss \parencite{JACOBS1985}. Iowa grass
waterways and drainageway alluvial soils reduced groundwater nitrate $70\%$
relative to cropped fields and removed $144$~kg~yr$^{-1}$
($7.8$~kg~ha$^{-1}$) of inflow nitrate \parencite{Schilling2022}. However,
where tile drains and ditches route water directly to streams they can
bypass the riparian zone, limiting attenuation \parencite{Puckett2005}.
Vegetated filter strips reduced surface-runoff nitrate, PO$_4$, and TP (up
to $70$--$75\%$ for P), but subsurface nitrate could increase through the
strip \parencite{Bhattarai2009}; regionally, filter strips were the most
effective practice for surface sediment, TN ($-57\%$), and TP ($-63\%$)
\parencite{Yuan2022}. Improved subsurface-drainage filter materials
(sand--gravel, zeolite) enhanced ammonium and phosphate removal and salt
discharge, though some configurations increased nitrate and TN
\parencite{Tao2019b,Tao2022,Tao2025b,Wang2024e}.

\subsubsection{Comparative effectiveness and cost}

Comparative syntheses provide load-reduction benchmarks. Denitrifying
bioreactors, constructed wetlands, cover crops, crop rotation, and nutrient
management each reduced subsurface nitrate loads by $23$--$40\%$
\parencite{Yuan2022}. In the Lower Mississippi Alluvial Valley, nine BMPs
delivered nutrient reductions of $15$--$100\%$, with mean TN reductions
highest for conservation tillage ($76\%$) and land-use conversion, and
subsurface drainage itself yielding TN and TP reductions of $17$ and $31\%$
\parencite{Kröger2012}. Cost-effectiveness analysis ranked wetlands most
cost-effective over both 10- and 50-year horizons ($2.04$--$2.91$~\$~kg$^{-1}$~N),
followed by two-stage ditches and cover crops, with USDA payments reducing
landowner costs $76$--$90\%$ \parencite{Roley2016}; wetland N-removal costs
across studies ranged $0.66$--$58$~\$~kg$^{-1}$~N \parencite{Messer2021}.
Farmer-facing evaluation found nutrient management most compatible and
wetlands richest in ecosystem services but least compatible, with cover
crops moderate and positive \parencite{Christianson2014}.

\subsection{Salinity, Trace Elements, Pesticides, and Contaminant Dynamics}

Beyond N and P, drainage governs salt, trace-element, and pesticide
transport. Reducing applied irrigation and the leaching fraction can cut
salt discharge by $50\%$ or more, and San Luis Drain data showed salt load
proportional to discharge at $9.2$~kg~m$^{-3}$; drainage-water reuse for
salt-tolerant crops and evaporation basins are practical pathways, though
basins accumulate extreme salinity ($4$--$390$~g~L$^{-1}$) and trace
elements \parencite{WESTCOT1988}. Selenium in San Joaquin drainage
motivated evaporation-pond hazard management, where only four of twenty
methods (removing levee vegetation and windbreaks, deepening ponds, hazing
birds) were recommended, and no single method could eliminate the hazard
\parencite{BRADFORD1991}; open-water evaporation is prohibited in
California owing to selenium, requiring rock-media systems
\parencite{Ayars2014}. Egyptian forecasts indicated pumped-groundwater TDS
would remain below $800$~ppm over 100 years with no crop water-quality
problems \parencite{RADSTAKE1988}, and Pakistani pipe-drainage salinity
showed essentially no temporal trend owing to lateral groundwater inflow
\parencite{Kelleners1998}. Israeli subsurface and bio-drainage systems
lowered water tables below $2$--$3$~m and improved the salinity--sodicity
regime on over $95\%$ of $3000$~ha \parencite{Gafni2001}. Wrapped
subsurface pipes with sand--gravel filters achieved salt-removal rates
above $95\%$ \parencite{Wang2024e}, and isotopic/hydrochemical analysis
distinguished irrigation, phreatic, and drainage water sources of salt in
saline-alkali land \parencite{Zhou2023}. Idaho irrigation return flow
carried elevated salts and $10$~kg~ha$^{-1}$~yr$^{-1}$ nitrate-N owing to
tunnel/tile flow, but concentrations declined significantly over
$2006$--$2018$ \parencite{Bjorneberg2020}.

Drainage of acid sulphate soils is a distinct hazard: subsurface drainage
was a major cause of stream pH declines to near or below $4.0$ (and $3.0$
in ditches), elevating Al and other metals, with controlled pipe drainage
recommended over in-stream liming \parencite{Osterholm2005,Åström2005}.
Pesticides appear in drainage within a month of application as a general
rule, with losses tied to application quantity and compound properties;
highest concentrations occur at plot scale and dilute downstream
\parencite{Tournebize2012}. Constructed wetlands and small grass-strip
wetlands remove pesticides with efficiency governed by sorption ($K_{oc}$),
degradation ($DT_{50}$), and residence time, sometimes with negative
efficiencies \parencite{Tournebize2017,Vallée2015}. Persistent degradates
such as metolachlor ESA remained in groundwater up to $18$~years but did
not migrate far, and no detections exceeded MCLs \parencite{Puckett2005};
drainage increased nitrate but decreased atrazine relative to runoff-only
conditions \parencite{SPALING1995}. Isotopic and geochemical tracers
resolve sources and pathways: combined stable isotopes and fecal markers
showed nitrate sources vary with hydrological stage and that soils play a
determinant role in nitrate formation and transfer \parencite{Briand2017},
and multi-tracer analysis distinguished shallow chloride- and nitrate-rich
groundwater from regional carbonate-aquifer water in the Wabash headwaters
\parencite{Frisbee2022}. Direct effluent bypass through mole-pipe drains
contributed disproportionately to pollutant loads (\emph{c.} $81\%$ of
\emph{E. coli}, $72\%$ of ammonium-N, $55\%$ of P) despite averaging only
$\sim0.7\%$ of applied effluent, motivating storage and slower application
\parencite{Monaghan2004,Monaghan2007}.

\subsection{Watershed-Scale Assessment, Modelling, and Reduction Strategies}

\subsubsection{Models and monitoring}

Process models are well validated for drainage hydrology and N. DRAINMOD
and its derivatives, RZWQM2, RZ-SHAW, and SWAT reproduce drainage flow,
water-table dynamics, and nitrate export across climates
\parencite{Youssef2018,Abdelbaki2025,Craft2018,Jiang2020,Tian2012,Mehan2019}.
LEACHM predicted nitrate leaching better under controlled
drainage/subirrigation than free drainage \parencite{Ng2000}, and
DRAINMOD-FOREST captured drainage and nitrate in managed pine plantations
\parencite{Tian2012}. Monitoring design strongly affects inferred loads:
high-frequency ($15$-min) sampling captured concentration ranges, diurnal
variation, and storm dynamics missed by grab sampling, which overestimated
loads \parencite{Drake2018}. Spatially distributed, event-based, and
year-round sampling is needed to resolve subwatershed heterogeneity and hot
spots \parencite{Irvine2019,Husk2024}, and long-term records are essential
because significant nitrate reductions are far more detectable in
flow-weighted concentration than in load or yield \parencite{Danalatos2022,Ghane2016b}.
Total-N concentrations decrease as spatial scale increases, implying
scale-specific water-quality criteria \parencite{Lagzdins2012}.

\subsubsection{Achieving reduction targets and residual challenges}

Watershed-planning studies show that meeting large reduction goals requires
combining in-field and edge-of-field practices. In the Mississippi basin,
nitrogen-management practices alone were insufficient (improved fertilizer
management and cover crops together cut delivered loads $30.2\%$), but a
full suite of N-management plus N-removal practices reached the $45\%$
target while converting under $6\%$, and as little as $2.5\%$, of cropland
when practices were combined and targeted \parencite{McLellan2015}.
High-resolution LiDAR-based planning identified scenarios achieving
$34$--$39\%$ nitrate reduction while removing $<5\%$ of cropland, with cover
crops and nutrient-removal wetlands common to the best scenarios
\parencite{Tomer2015}. Long-term simulation suggested DWM could plausibly
help meet Iowa's $41\%$ reduction goal \parencite{Craft2018}, and a Czech
pilot designed measures reducing drainage nitrate load $24\%$ across $38\%$
of the drained area \parencite{Kvítek2023}. Yet watershed-scale evidence
also cautions that traditional surface-water BMPs may be inadequate over
tile drainage: seven years of increased strip-till, grassed waterways, and
stream buffers produced no measurable change in nitrate, phosphorus,
sediment, or hydrology, because they did not intercept subsurface flow,
which must be a design priority \parencite{Lemke2011}. Ecosystem-service
framing shows drainage modernization can sharply reduce migratory-bird
habitat (up to $87\%$ for dabbling ducks) while paired water-quality
wetlands substantially reduce nitrate export and partially offset habitat
loss, benefiting diving ducks most \parencite{Mitchell2022,Mitchell2025}.
Managed and semi-natural ditch networks themselves provide biodiversity,
water-purification, and nutrient-retention functions that depend on
vegetation, morphology, and maintenance, with no single maintenance regime
optimizing all services \parencite{Dollinger2015,Herzon2008}. Field
observations that subsurface drainage installation did not measurably alter
soil-health metrics or soil organic carbon over two years, against strong
seasonal and interannual variation, further temper expectations of rapid
in-soil change \parencite{Sherbine2023,Huang2023}.